\documentclass[11pt]{article}
\usepackage[a4paper,margin=1in]{geometry}
\usepackage{fontspec}
\usepackage[boldfont=false]{xeCJK}
\setCJKmainfont{FandolSong-Regular.otf}
\usepackage{amsmath}
\usepackage{amssymb}
\usepackage{array}
\usepackage{booktabs}
\usepackage{graphicx}
\usepackage{adjustbox}
\usepackage{enumitem}
\setlist[itemize]{topsep=2pt,partopsep=0pt,parsep=0pt,itemsep=2pt,leftmargin=1.4em}
\setlist[enumerate]{topsep=2pt,partopsep=0pt,parsep=0pt,itemsep=2pt,leftmargin=1.6em}
\usepackage{parskip}
\usepackage{fvextra}
\fvset{breaklines=true,breakanywhere=true,fontsize=\small}
\setlength{\parindent}{0pt}

\begin{document}

\begin{center}
  {\LARGE\bfseries Halogen compounds}\\[0.35em]
  {\large A Level Chemistry}
\end{center}
\vspace{0.6em}

\section*{Halogenoalkanes}

A \textbf{halogenoalkane} 卤代烷 is an alkane with one or more halogen atoms in place of hydrogen (a C–X bond, where X is a halogen).

\subsection*{Making halogenoalkanes}

\begin{itemize}
\item \textbf{free-radical substitution} 自由基取代 of an alkane with $\text{Cl}_2$ or $\text{Br}_2$ in ultraviolet light.

\item \textbf{electrophilic addition} 亲电加成 of an alkene with a halogen $\text{X}_2$ or a hydrogen halide $\text{HX}$.

\item substitution of an \textbf{alcohol} 醇, for example by $\text{HX}$, by $\text{PCl}_5$, by $\text{PCl}_3$ with heat, or by $\text{SOCl}_2$.

\end{itemize}

\subsection*{Three classes}

A halogenoalkane is \textbf{primary} 伯, \textbf{secondary} 仲 or \textbf{tertiary} 叔, depending on how many carbon atoms are joined to the carbon that holds the halogen (one, two or three).

\section*{Nucleophilic substitution}

The C–X bond is polar, so the carbon is slightly positive. A \textbf{nucleophilic substitution} 亲核取代 happens when a \textbf{nucleophile} 亲核试剂 (a lone-pair species) attacks that carbon and replaces the halogen.

\begin{center}
\adjustbox{max width=\linewidth}{%
\begin{tabular}{ll}
\toprule
\textbf{Reagent and conditions} & \textbf{Product} \\
\midrule
$\text{NaOH(aq)}$, heat & an alcohol \\
$\text{KCN}$ in ethanol, heat & a \textbf{nitrile} 腈 (adds one carbon to the chain) \\
$\text{NH}_3$ in ethanol, heated under pressure & an \textbf{amine} 胺 \\
\bottomrule
\end{tabular}}
\end{center}

To \textbf{identify the halogen}, warm the halogenoalkane with \textbf{silver nitrate} 硝酸银 in ethanol. A silver halide \textbf{precipitate} 沉淀 forms, and its colour shows which halogen is present (white $\text{AgCl}$, cream $\text{AgBr}$, yellow $\text{AgI}$).

\section*{Elimination}

The same halogenoalkane can instead undergo \textbf{elimination} 消去 to form an \textbf{alkene} 烯烃. The conditions decide which reaction wins:

\begin{itemize}
\item $\text{NaOH}$ in \textbf{water} → nucleophilic substitution → an alcohol.

\item $\text{NaOH}$ in \textbf{ethanol}, heated → elimination → an alkene.

\end{itemize}

\[
\text{C}_2\text{H}_5\text{Br} + \text{NaOH} \rightarrow \text{C}_2\text{H}_4 + \text{NaBr} + \text{H}_2\text{O}
\]

\section*{The S$_\text{N}$1 and S$_\text{N}$2 mechanisms}

Nucleophilic substitution can follow two routes:

\begin{itemize}
\item \textbf{S$_\text{N}$2}: one step. The nucleophile attacks at the same time as the halogen leaves, passing through a crowded \textbf{transition state} 过渡态 where both are half-bonded. The rate depends on both the halogenoalkane and the nucleophile.

\item \textbf{S$_\text{N}$1}: two steps. First the C–X bond breaks to give a \textbf{carbocation} 碳正离子; then the nucleophile attacks it. The rate depends only on the halogenoalkane.

\end{itemize}

Alkyl groups push electrons towards the positive carbon (the \textbf{inductive effect} 诱导效应), so they stabilise the carbocation. This is why:

\begin{itemize}
\item \textbf{primary} halogenoalkanes mostly react by S$_\text{N}$2.

\item \textbf{tertiary} halogenoalkanes mostly react by S$_\text{N}$1 (their carbocation is well stabilised).

\item \textbf{secondary} halogenoalkanes use a mixture of the two.

\end{itemize}

\section*{Different reactivities}

How fast a halogenoalkane reacts depends on the strength of the C–X bond, measured by its \textbf{bond energy} 键能. The C–I bond is the weakest, so iodoalkanes react fastest; the C–Cl bond is the strongest of the three, so chloroalkanes react slowest. So when tested with silver nitrate, an iodoalkane gives its precipitate first — its higher \textbf{reactivity} 反应活性 comes from the weaker C–X bond.


\end{document}
